\section{2. Related Work}

\subsection{Traditional Light Field Depth Estimation Methods}

Many depth estimation methods for light field cameras have been proposed based on Epipolar Plane Images (EPIs) [9]. These approaches explicitly construct EPIs and derive depth by calculating the slopes of linear structures, relying heavily on the photo-consistency constraint across different viewpoints. For instance, Wanner et al. [10] proposed a globally optimized framework using structure tensors to accurately extract EPI slopes, achieving notable precision on Lambertian surfaces. Similarly, recent work [11] has extensively explored edge-aware filtering and locally polynomial fitting to refine slope estimation in texture-rich regions. However, most of them rely on the Lambertian assumption and work poorly on glossy surfaces, where the angular photo-consistency is substantially violated by viewpoint-dependent reflectance.

While EPI-based methods primarily exploit angular consistency, defocus-based approaches extend the geometric reasoning to the spatial domain by analyzing focal stacks [12]. These methods evaluate the sharpness or focus metrics across a stack of synthetically refocused images to determine the optimal depth. Tao et al. [13] successfully combined defocus and correspondence cues in a unified framework, significantly improving depth accuracy in occluded regions. Additionally, Lin et al. [4] introduced an adaptive focus metric to properly handle spatially-varying (SV) textures. Nonetheless, these focus-based metrics inherently assume that the radiance of a 3D point remains invariant across different viewpoints. When dealing with non-Lambertian materials, the viewpoint-dependent reflectance introduces aliasing bimodal distributions in the focal stack, causing the focus metrics to fail and resulting in severe depth artifacts.

To mitigate the failures of pure geometric cues on complex surfaces, traditional methods for non-Lambertian surfaces incorporate explicit Bidirectional Reflectance Distribution Function (BRDF) modeling or photometric stereo techniques [14]. Shape from shading and photometric stereo have a long history. Since these are highly under-constrained problems, most work assumes a known light source to increase feasibility. For example, Barron and Malik [15] described a framework that utilizes the dichromatic reflection model to separate diffuse and specular components, thereby recovering both shape and reflectance. A follow-up work [16] adopts this model to eliminate specular pixels before depth estimation. However, the dichromatic model fails to hold for materials like metals or complex anisotropic scatterers. Furthermore, none of these methods tries to recover shape or reflectance using passive camera motions, and the techniques are not intended for light field cameras, as they typically require active illumination or highly complex priors.

In short, traditional light field depth estimation methods exhibit critical limitations when confronting real-world complex scenes. First, they severely rely on the pure Lambertian assumption; in non-Lambertian regions such as transparent or mirror (ToM) materials and strong scatterers, the violation of viewpoint consistency causes EPI slope calculations to fail completely, degrading depth prediction into mere texture copying. Second, traditional approaches designed for non-Lambertian surfaces (e.g., photometric stereo) typically necessitate active light sources or extremely complex priors, making them difficult to directly apply to passive light field acquisition scenarios. Third, these methods lack pixel-level adaptive perception capabilities for mixed materials (coexistence of Lambertian and non-Lambertian), leading to severe artifacts in depth maps under complex scenes and failing to provide reliable physical priors for modern deep learning frameworks. In contrast, our proposed unified dual-mask physical model explicitly separates and perceives material components from a physical optics perspective. This provides robust physical priors for mixed scenes and substantially resolves the depth ambiguity in non-Lambertian regions, addressing the critical limitations of previous light field shape acquisition approaches.

\subsection{Deep Learning-based General Light Field Depth Estimation Methods}

While traditional methods explicitly formulate geometric constraints, deep learning-based approaches have recently dominated light field depth estimation by learning implicit mappings from large-scale datasets [17]. These methods primarily utilize Convolutional Neural Networks (CNNs) or Transformers to extract geometric and textural features from various light field representations, such as Epipolar Plane Images (EPIs), macro-pixels, and multi-view sub-aperture images [18]. By leveraging massive amounts of annotated data, these data-driven paradigms have significantly pushed the boundaries of depth estimation accuracy in controlled, Lambertian-dominant environments.

Pioneering this direction, Shin et al. [19] proposed EPINET, which explicitly separates spatial and angular information. By feeding a central view and its corresponding angular views into a multi-stream CNN [20], EPINET successfully fuses spatial textures with angular disparities to infer depth. Although this architecture significantly improves accuracy on standard benchmarks, it implicitly relies on the photo-consistency assumption during the angular feature fusion stage. Consequently, the network struggles to resolve depth ambiguities in non-Lambertian regions where viewpoint-dependent reflectance violates this fundamental assumption.

To better capture the intrinsic 4D structure of light fields, subsequent works have explored 4D convolutions and spatial-angular disentangled modules to separate features more effectively [21]. For instance, Zhou et al. [22] introduced a pseudo-4D convolution to efficiently model spatial and angular correlations without incurring prohibitive computational costs. Similarly, Wang et al. [23] developed a disentangled network to explicitly separate spatial and angular features, significantly enhancing depth estimation in occluded and texture-less regions. While these methods dramatically improve the representation capacity for complex geometric structures, they substantially treat spatially-varying (SV) non-Lambertian reflections as unstructured noise. Without explicit physical priors to separate specular components from diffuse albedo, the learned features remain heavily entangled in mixed-material scenes, violating the working principles of pure geometric matching.

More recently, Transformer-based architectures have been adapted to light field depth estimation to capture long-range dependencies that CNNs typically miss [24]. Methods such as the Light Field Transformer [25] and Epipolar Transformer [26] leverage self-attention mechanisms to model global spatial-angular correlations. While these Transformer-based models significantly enhance depth accuracy on standard benchmarks by capturing broader contextual information, they still implicitly assume photo-consistency and substantially struggle with materials like metals or complex anisotropic scatterers where angular appearances vary drastically.

Furthermore, some approaches reformulate the light field as pseudo-sequences or multi-view stereo setups [27], employing recurrent neural networks or cost volume aggregation to infer depth. Although this strategy significantly improves depth accuracy on global structures, it often fails to preserve fine-grained details in non-Lambertian regions due to the severe violation of epipolar geometry.

In short, despite their considerable empirical success, existing deep learning-based methods predominantly rely on the Lambertian assumption. When confronted with non-Lambertian materials, the viewpoint-dependent reflectance severely disrupts the learned geometric features. While some recent works attempt to incorporate physical priors or regularize the feature distribution, substantially resolving the failure of traditional geometric matching and the texture-copying issue to some extent, they lack a unified physical model to explicitly decouple mixed materials. Consequently, there remains a critical need for a physically-grounded framework that can intrinsically guide deep networks to handle complex, real-world light field scenes.